\documentclass[12pt,letterpaper]{article}
\usepackage{amsmath}
\usepackage{amsthm}
\usepackage{amsfonts}
\usepackage{amssymb}
\usepackage{amscd}
\usepackage{enumerate}
\usepackage{fancyhdr}
\usepackage{mathrsfs}
\usepackage{bbm}
\usepackage{framed}
\usepackage{mdframed}
\usepackage{cancel}
\usepackage{mathtools}
\usepackage{verbatim}
\usepackage{hyperref}
\usepackage[letterpaper,voffset=-.5in,bmargin=3cm,footskip=1cm]{geometry}
\setlength{\parindent}{0.0in}
\setlength{\parskip}{0.1in}
\allowdisplaybreaks
\headheight 15pt
\headsep 10pt
\newcommand\N{\mathbb N}
\newcommand\Z{\mathbb Z}
\newcommand\R{\mathbb R}
\newcommand\Q{\mathbb Q}
\newcommand\lcm{\operatorname{lcm}}
\newcommand\setbuilder[2]{\ensuremath{\left\{#1\;\middle|\;#2\right\}}}
\newcommand\E{\operatorname{E}}
\newcommand\V{\operatorname{V}}
\newcommand\Pow{\ensuremath{\operatorname{\mathcal{P}}}}

\DeclarePairedDelimiter\ceil{\lceil}{\rceil}
\DeclarePairedDelimiter\floor{\lfloor}{\rfloor}
\newcommand\hint[1]{\textbf{Hint}: #1}
\newcommand\note[1]{\textbf{Note}: #1}
\newcounter{enum22i}
\newenvironment{22enumerate}{\begin{enumerate}[a.]\itemsep0em\setcounter{enumi}{\value{enum22i}}}{\setcounter{enum22i}{\value{enumi}}\end{enumerate}}
\newenvironment{22itemize}{\begin{itemize}\itemsep0em}{\end{itemize}}
\fancypagestyle{firstpagestyle} {
  \renewcommand{\headrulewidth}{0pt}
  \lhead{\textbf{CSCI 0220}}
  \chead{\textbf{Discrete Structures and Probability}}
  \rhead{M. Littman}
}

\fancypagestyle{fancyplain} {
  \renewcommand{\headrulewidth}{0pt}
  \lhead{\textbf{CSCI 0220}}
  \chead{Problem}
  \rhead{\textit{ }}
}
\pagestyle{fancyplain}
\usepackage{tikz}

\begin{document}
  \thispagestyle{firstpagestyle}
  \begin{center}
    {\huge \textbf{Homework 3}}

    {\large Due:\textit{ Monday, June 7th at 11:59pm EDT}}
  \end{center}


     All homeworks are due the Monday after their release at 11:59pm EDT on Gradescope.

    Please do not include any identifying information about yourself in the handin, including your Banner ID.

    Be sure to fully explain your reasoning and show all work for full credit. We recommend checking out the sample proofs on the course website to see the level of rigor for which we're looking.
  
  
  

\subsection*{Problem 1}
{
\nopagebreak 

{\setcounter{enum22i}{0}
      For this problem, let
	    \begin{align*}
    		A &= \{1,2,3,4,5\}\\
    		B &= \Pow(\{1,4,9\})\\
    		C &= \{x \mid \exists \ y \in \Z \text{ such that } y^2=x, \ x < 10 \}
	    \end{align*}
      \begin{22enumerate}
    		\item Find the following sets:
      		\begin{enumerate}[i.]
          		\item 
          	$A \cup C$
          		\item 
          	$A \cap B$
                \item
              $A - C$
      			\item
              $C \times C$
      			\item
              $\{x \mid x \in C \text{ AND } \exists\;y,z \in A \text{ such that } y+z = x\}$
    		  \end{enumerate}
          \item Find the cardinalities of the following sets:
            \begin{enumerate}[i.]
              \item $C \times C$
              \item $\Pow(C \times C)$
              \item $\Pow(\Pow(A) \cup B)$
            \end{enumerate}
            
    \item Assume that the sets $A, B, C$ are all
    subsets of the same universal set $U$. \textbf{Using set algebra}, show that the following equivalences hold.
	\begin{enumerate}[i.]
		\item $(A - B)\cup (B - A) 
        = (A\cup B) - (A\cap B)$
		\item $A \cup (B \cap (C - A)) = A \cup (B \cap C)$
        \item $\overline{(A \cup (B \cap C))} = (\overline{C} \cup \overline{B} ) \cap \overline{A}$
	\end{enumerate}
      \end{22enumerate}
}
}


\subsection*{Problem 2}
{
\nopagebreak 

{\setcounter{enum22i}{0}

        Consider a finite set $S$ and its power set, $ \Pow (S)$. For parts a through d, find and prove the size of the relevant set. If it is non-empty, state which elements of $\Pow(S)$ are in it and explain why.
        \begin{enumerate}[a.]
            \item $K_1=\{x \in \Pow (S) \; | \; \forall y \in \Pow (S), \; x \cap y = \emptyset \}$.
            \item $K_2=\{x \in \Pow (S) \; | \; \forall y \in \Pow (S), \; x \cap y = x \}$.
            \item $K_3=\{x \in \Pow (S) \; | \; \forall y \in \Pow (S), \;  x \cap y = y \}$.
            \item $K_4=\{x \in \Pow (S) \; | \; \forall y \in \Pow (S), \; x \cup y = x \}$.
            \item What is the intersection of $K_1$, $K_2$, $K_3$, and $K_4$?
            \item What is the union of $K_1$, $K_2$, $K_3$, and $K_4$?
        \end{enumerate}
}
}


\subsection*{Problem 3}
{
\nopagebreak 

{\setcounter{enum22i}{0}
A binary operation, $\star$, is an operation on a set that takes in two elements from the set and returns a third. For example, the addition operation $+$ is a binary operation on the integers.

We say that a binary operation $\star$ on a set $S$ is \textit{commutative}  if \[x\star y=y\star x \, \textrm{ for all } \, x,y\in S.\]

Let $X$ be a finite set. For each of the following operations, prove whether or not the operation is commutative over $\Pow(X)$.
\begin{enumerate}[i.]
	\item Set union
	\item Set intersection
	\item Set difference
	\item Symmetric difference
\end{enumerate}

%  (Symmetric difference can be defined in several equivalent ways including $(A \cup B)-(A \cap B)$.) % See the lecture notes and/or text for definitions of any other needed operations.)
}
}

\subsection*{Mind Bender (Extra Credit)}
{
\nopagebreak 

{\setcounter{enum22i}{0}
Mad about the chaos that ensued from the game of asteroid gravitational deflection, Venus is looking to instead host the parties for all future syzygies. It looks to invite distinct groups of celestial bodies to each one. To do so, it begins writing out the powerset of the set of celestial bodies in the solar system.\\
\\
Bored of how surprisingly long it's taking, Venus starts to think about different properties of the powerset. It wonders: if the powerset of the powerset of the set A consists of A, the powerset of A, A’s only element, and A’s only element’s only element, then what is A? Help Venus answer this question to get back to planning.\\
\\
\textbf{Note: }Make sure to prove your answer is correct \emph{and} unique.

}
}
\end{document}